\section{The integral decomposition complex}\label{sec:integral-decomposition}

Let \(i:\{0\}\hookrightarrow\Theta\).  Proper base change and Thom
excision identify
\begin{equation}\label{eq:stalk-costalk}
 i^*\mathcal K\cong R\Gamma(X,\Z)[4],\qquad
 i^!\mathcal K\cong R\Gamma(X,\Z)[2].
\end{equation}
Under these identifications, the natural map from costalk to stalk is cup
product by the Euler class \(-h\).  On the algebraic cohomology of \(X\),
its three nonzero blocks are
\begin{equation}\label{eq:three-intersection-blocks}
 H^0(X)\xrightarrow{-h}H^2(X),\qquad
 H^2(X)\xrightarrow{-h}H^4(X),\qquad
 H^4(X)\xrightarrow{-h}H^6(X),
\end{equation}
with matrices \([-1],[-3],[-1]\) in the generators
\(1,h,\ell,[\mathrm{pt}]\).  The cone is
\(R\Gamma(K,\Z)[4]\), so this is also the Gysin triangle behind
Lemma~\ref{lem:link-cohomology}.

We use a standard intersection-form description of point summands.  By
adjunction, morphisms from and to \(\Z_0[m]\) are Borel--Moore classes of
the fibre, and their composition is the corresponding fibre intersection
form.  De Cataldo and Migliorini develop the rational mechanism for
projective resolutions, including isolated fourfold singularities
\cite[Sections~2.4 and~4.4]{deCataldoMigliorini}.  Over a field the modular
rank criterion is
\cite[Proposition~3.2 and Corollary~3.5]{JuteauMautnerWilliamson}; the same
adjunction calculation over \(\Z\) applies when the displayed rank-one form
is a unit.  The first and third blocks of
\eqref{eq:three-intersection-blocks} therefore give inclusion--projection
pairs with identity composite.  Splitting their idempotents yields
\[
 \mathcal K\cong\Z_0[2]\oplus\mathcal P\oplus\Z_0[-2].
\]
The cross-compositions vanish for degree reasons.  The remaining object has
stalk cohomology only in degrees \(-4,-1,0\), costalk cohomology only in
degrees \(0,1,4\), and restriction \(\Z_U[4]\).  These are the support and
cosupport conditions for \(\mathcal P\) to be perverse.
All its displayed stalk and costalk groups are free.  Consequently derived
reduction modulo three preserves these degree bounds and hence perversity;
the same freeness excludes a torsion point summand of \(\mathcal P\).

Its central attachment is the exact sequence
\begin{equation}\label{eq:central-zigzag}
 H^3(K,\Z)\longrightarrow
 A=\Z h\xrightarrow{-3}B=\Z\ell
 \longrightarrow H^4(K,\Z).
\end{equation}
The first arrow is zero, while the last sends \(\ell\) to
\(\tau=p^*\ell\), the order-three class in the link.  Indeed,
\eqref{eq:central-zigzag} is the degree-three/four part of the Gysin
sequence.  This proves the integral decomposition and attachment asserted
in Theorem~\ref{thm:integral-decomposition}.

Over \(\Z[1/3]\), the middle map in
\eqref{eq:central-zigzag} is a unit.  Its normalized inclusion--projection
pair splits a central point sheaf.  The complement extends \(\Z_U[4]\) and
has neither a point subobject nor a point quotient, hence is the intermediate
extension.  Thus
\begin{equation}\label{eq:localized-decomposition}
 R\sigma_*\Z[1/3]_M[4]\cong
 \IC_\Theta(\Z[1/3])\oplus
 \Z[1/3]_0[2]\oplus\Z[1/3]_0\oplus\Z[1/3]_0[-2].
\end{equation}
The same proof works over every field of characteristic different from
three.  Integrally, a free central point summand would contribute a unit block
\(\Z\to\Z\).  Since \(A\) and \(B\) both have rank one, this contradicts
the Smith factor three in \eqref{eq:central-zigzag}.

In characteristic three, the universal-coefficient sequence and
Lemma~\ref{lem:link-cohomology} give
\[
 \dim_{\F_3}H^3(K,\F_3)=
 \dim_{\F_3}H^4(K,\F_3)=11.
\]
The derived reduction of \eqref{eq:central-zigzag} is
\begin{equation}\label{eq:mod-three-zigzag}
 \F_3^{11}\twoheadrightarrow\F_3
 \xrightarrow{0}\F_3\hookrightarrow\F_3^{11}.
\end{equation}
Write \(\mathcal P_3=\mathcal P\otimes^{L}_{\Z}\F_3\); the freeness noted
above shows directly that \(\mathcal P_3\) is perverse.
The extra degree-three class is the Tor, or Bockstein, companion of
\(\tau\).  The zero middle form rules out a point summand.  Any direct-sum
decomposition of \(\mathcal P_3\) has at most one summand with
nonzero restriction to \(U\), because that restriction is the simple
rank-one local system; every other summand would be point-supported.
Therefore the residual object is indecomposable.  Its nonzero kernel and
cokernel in \eqref{eq:mod-three-zigzag} show that it is not the intermediate
extension of that simple local system, hence is not simple.  An
indecomposable semisimple object is simple, so the residual object is
nonsemisimple.
This direct argument is the point-stratum case of Cipriani's general
small-extension criterion
\cite[Lemma~3.12, Theorem~3.21, and Corollary~3.22]{Cipriani}; the link
cohomology, Smith factor three, and the two dimensions \(11\) are the
example-specific inputs here.

The same criterion gives the full Loewy structure.  Put
\(\delta=\F_{3,0}\), and write
\(A={}^pH^0(i^!\mathcal P_3)\) and
\(B={}^pH^0(i^*\mathcal P_3)\).  Equation
\eqref{eq:mod-three-zigzag} says that \(A=B=\F_3\) and that the canonical
map \(A\to B\) is zero.  Cipriani's projection sequences therefore give a
canonical filtration
\[
 0\subset\delta\subset\mathcal Q
   \subset\mathcal P_3
\]
whose successive quotients are
\(\delta\), \(\IC_\Theta(\F_3)\), and \(\delta\).  The first adjacent
extension cannot split because \(\mathcal Q\) has no point quotient; the
second cannot split because \(\mathcal P_3/\delta\) has no point subobject.
The canonical projection sequences make this the unique composition series.
Thus \(\mathcal P_3\) is uniserial of length three, completing the proof of
Theorem~\ref{thm:integral-decomposition}.

\begin{proof}[Proof of Corollary~\ref{cor:modular-rhl}]
The class \(\eta=c_1(\mathcal O_M(-X))\) restricts to \(h\) on the
exceptional cubic.  Under the unit-normalized outer splittings above,
\({}^pH^{-2}(\mathcal K)\) and \({}^pH^2(\mathcal K)\) correspond on the
central stalk to the generators \(h\in H^2(X,\Z)\) and
\([\mathrm{pt}]\in H^6(X,\Z)\), respectively.  Therefore
\[
 \eta^2(h)=h^3=3[\mathrm{pt}].
\]
The relative Lefschetz map is multiplication by three, so its reduction
modulo three is zero rather than an isomorphism.
\end{proof}

We finish by recording the global fate of \(\tau\).  Apply
Lemma~\ref{lem:clean-base-change} to
\(T=[\mathrm{pt}]\otimes1\in H^4(F\times F,\Z)\), and put
\[
 u_4=q_*\mu^*T\in H^4(M,\Z).
\]
Then
\begin{equation}\label{eq:degree-four-fano-lift}
 e^*u_4=\ell,\qquad b_*u_4=[F]=\theta^{[3]}.
\end{equation}
The restriction of \(u_4|_U\) to the link is \(\tau\).  This class has
infinite order: if it vanished rationally, then \(u_4\) would be rationally
supported on the contracted divisor and would have zero pushforward to
\(J\), contrary to \eqref{eq:degree-four-fano-lift}.

Rationally, \(H^4(U)\to H^4(K)\) is zero.  One way to see this is to combine
\eqref{eq:kraemer-decomposition} with the pair sequence for \((M,U)\): the
exceptional map \(H^2(X,\Q)\to H^4(M,\Q)\) is injective, and
\(H^3(X,\Q)\to H^5(M,\Q)\) is injective by
Theorem~\ref{thm:main-escape}.  Hence \(H^4(U,\Q)\) and
\(IH^4(\Theta,\Q)\) have the same dimension, while the truncation map from
the latter to the former is injective.  Thus the integral boundary image is
torsion, and \eqref{eq:degree-four-fano-lift} shows that it is exactly
\(\langle\tau\rangle\).  The truncation triangle now gives
\begin{equation}\label{eq:degree-four-index-three}
 0\longrightarrow IH^4(\Theta,\Z)\longrightarrow H^4(U,\Z)
 \longrightarrow\Z/3\longrightarrow0.
\end{equation}

Finally, normality and the Deligne model give a triangle
\[
 \Z_\Theta[4]\longrightarrow\IC_\Theta(\Z)
 \longrightarrow H^3(K,\Z)_0[1]\xrightarrow{+1}.
\]
The map \(IH^3(\Theta,\Z)=H^3(U,\Z)\to H^3(K,\Z)\) is surjective by
Propositions~\ref{prop:topological-exact} and
\ref{prop:endpoint-lifts}.  The connecting map to \(H^4(\Theta,\Z)\)
therefore vanishes.  It follows that
\begin{equation}\label{eq:ordinary-equals-intersection-above-middle}
 H^k(\Theta,\Z)\xrightarrow{\sim}IH^k(\Theta,\Z)
 \qquad(k\geq4).
\end{equation}
Together, \eqref{eq:degree-four-index-three} and
\eqref{eq:ordinary-equals-intersection-above-middle} show that the local
order-three class is an index-three attachment, not a global torsion class.
